% =========================================================
% frontmatter/objective.tex
% =========================================================

\chapter*{Project Objective}
\thispagestyle{fancy}
\addcontentsline{toc}{chapter}{Project Objective}
\begingroup
\setlength{\emergencystretch}{2em}

The objective is to establish and verify the modal behavior of an ideal hollow waveguide with a square \(20~\mathrm{mm}\times20~\mathrm{mm}\) cross section and a physical length of \(90~\mathrm{mm}\). The study focuses on operation at \(10~\mathrm{GHz}\) over an \(8\)--\(15~\mathrm{GHz}\) simulation band.

The work addresses five engineering questions:

\begin{enumerate}
  \item Which low-order TE and TM modes are above cutoff at \(10~\mathrm{GHz}\)?
  \item How closely do CST port-eigenmode cutoff frequencies agree with closed-form rectangular-waveguide theory?
  \item How does a square cross section change the meaning of same-index and cross-index modal S-parameters?
  \item Do the electric-field, magnetic-field, and wall-current distributions support the predicted propagating and near-cutoff behavior?
  \item What limitations must be considered before using this idealized section as a practical horn feed?
\end{enumerate}

The evidence is analytical and simulated. The model uses vacuum filling and ideal PEC walls; finite conductor loss, fabrication tolerances, transitions, and measured validation are outside the present scope.
\par\endgroup
